% GENERATED FILE. Edit canonical lessons, then run npm run generate:books.
% canonical-source-hash: fnv1a64:ce422bed4030868a
% canonical-lessons: ZH-C364-quan, ZH-C364-buzhun, ZH-W364-zhun-allow, ZH-C364-biaozhun, ZH-C364-pizhun, ZH-C364-tixing

\chapter{Not Allowed}
\label{ch:zh-not-allowed}
\hlchaptermodality{364}

\begin{chapteropening}
\textbf{By the end of this chapter:} \emph{I can say what is not allowed, urge and remind someone, and say a plan was approved.}

Complete the last lesson of chapter 364: I can say what is not allowed, urge and remind someone, and say a plan was approved.
\end{chapteropening}

\section[quàn]{quàn — to urge; to try to persuade}

\label{lesson:ZH-C364-quan}



\begin{quote}
Before the new one: say the Mandarin for to break the law, then the Mandarin for a person's character and upbringing.
\end{quote}

\subsection*{You'll want to know: quàn}
\emph{quàn} — \textquotedblleft{}to urge; to try to persuade\textquotedblright{}.

\begin{quote}
\emph{wǒ quàn tā zuìhǎo bié hējiǔ} — I urged him not to drink.
\end{quote}

Its characters come later; for now it is a word for the ear.

\begin{cousinweb}
\emph{quàn} is inherited from Middle Chinese \emph{khjwonH} (Guangyun, in Baxter's transcription), which glosses it 'to encourage' and 'to teach'.
\end{cousinweb}

\begin{grammarlens}[title={What you've built}]
One more word for giving reasons.
\end{grammarlens}

\subsection*{Your turn}
\begin{itemize}
\raggedright
  \item \emph{Say it:} \emph{quàn}
  \item \emph{Say it:} \emph{quàn}, once more
  \item \emph{Say it:} \emph{wǒ quàn tā zuìhǎo bié hējiǔ} — I urged him not to drink.
  \item \emph{Recall:} say the Mandarin for to break the law, then the Mandarin for a person's character and upbringing, then say \emph{quàn} again
\end{itemize}

\begin{tcolorbox}[breakable,colback=teal!4,colframe=teal!35!black,title={Before you move on}]
What does quàn mean? (\textquotedblleft{}To urge; to try to persuade\textquotedblright{}.) Say it once more.
\end{tcolorbox}
\section[\texorpdfstring{\zh{不准} (bùzhǔn)}{不准 (bùzhǔn)}]{\zh{不准} (bùzhǔn) — to not allow}

\label{lesson:ZH-C364-buzhun}



\begin{quote}
Before the new one: say the Mandarin for a person's character and upbringing, then the Mandarin for to urge someone.
\end{quote}

\subsection*{You'll want to know: \zh{不准}}
\textbf{\zh{不准}} — \emph{bùzhǔn} — \textquotedblleft{}to not allow\textquotedblright{}.

\begin{quote}
\textbf{\zh{上课的时候不准吃东西}} \emph{shàngkè de shíhou bùzhǔn chī dōngxi} — No eating in class.
\end{quote}

The new character, in pieces: \zh{准} — ten strokes in two pieces: on the left the ice piece of \zh{冷}, two strokes; on the right a bird piece, eight strokes: a left-falling stroke, a long vertical, a dot, three short bars, a vertical down through them, and a long bar at the foot. The next lesson puts it in your hand.

The tone of bù stays falling here, because the next syllable is not a falling tone.

\begin{cousinweb}
\textbf{\zh{不准}} (\emph{bùzhǔn}) is built from \zh{不} 'not' and \emph{zhǔn}, both inherited from Middle Chinese, \emph{pjut} and \emph{tsywinX} (Guangyun, in Baxter's transcription), which lists the second as the everyday spelling of a word glossed 'even' and 'level'. How it came to mean allowing is not settled.
\end{cousinweb}

\begin{grammarlens}[title={What you've built}]
One more word for giving reasons.
\end{grammarlens}

\subsection*{Your turn}
\begin{itemize}
\raggedright
  \item \emph{Say it:} \emph{bùzhǔn}
  \item \emph{Say it:} \emph{bùzhǔn}, once more
  \item \emph{Say it:} \emph{shàngkè de shíhou bùzhǔn chī dōngxi} — No eating in class.
  \item \emph{Recall:} say the Mandarin for a person's character and upbringing, then the Mandarin for to urge someone, then say \emph{bùzhǔn} again
\end{itemize}

\begin{tcolorbox}[breakable,colback=teal!4,colframe=teal!35!black,title={Before you move on}]
What does \zh{不准} mean? (\textquotedblleft{}To not allow\textquotedblright{}.) Say it once more.
\end{tcolorbox}
\section[\texorpdfstring{\zh{准} (zhǔn)}{准 (zhǔn)}]{\zh{准} — the ice piece, then the bird piece}

\label{lesson:ZH-W364-zhun-allow}



\begin{quote}
The word first, then the character inside it.

\begin{itemize}
\raggedright
  \item \emph{Recall:} say \emph{bùzhǔn} — \textbf{R1}, one lesson back, before it decays
\end{itemize}

\begin{quote}
\zh{冷} \zh{不}
\end{quote}

\emph{Read it:} each one aloud, with nothing said first — \emph{lěng}, \emph{bù}
\end{quote}

\subsection*{Script — the ice piece, then the bird piece}
\hlblockfigure{\detokenize{figures/ZH-W364-zhun-allow-filmstrip.pdf}}{How \zh{准} is written, stroke by stroke}

\begin{quote}
\zh{准}
\end{quote}

\textbf{Two pieces, left before right.} Count the pieces, not the strokes.

\textbf{The left piece, the ice piece, two strokes}, as in \zh{冷}: a dot, then a stroke rising from the bottom left.

\textbf{The right piece, the bird piece, eight strokes}: a left-falling stroke, a long vertical down the left side, a dot at the top right, three short horizontals, a vertical down through them, then a long horizontal at the foot.

2 + 8 = ten.

\subsection*{Your turn}
\begin{itemize}
\raggedright
  \item \emph{Write it:} \zh{准}, all ten, with the model in front of you
  \item \emph{Cover:} write it again from memory, and count the pieces, not the strokes
\end{itemize}

\begin{tcolorbox}[breakable,colback=teal!4,colframe=teal!35!black,title={Before you move on}]
How many pieces in \zh{准}? (\textbf{Two}.) Which character you know has the same left piece? (\textbf{\zh{冷}}.) What closes the bird piece? (\textbf{A long horizontal at the foot}.)

Source: \href{https://raw.githubusercontent.com/chanind/hanzi-writer-data/68d10a4b21150cae5e1ebbd223eed289cf32d90c/data/\%E5\%87\%86.json}{Hanzi Writer Data \zh{准}, pinned stroke paths and medians}.
\end{tcolorbox}
\section[biāozhǔn]{biāozhǔn — standard to measure by}

\label{lesson:ZH-C364-biaozhun}



\begin{quote}
Before the new one: say the Mandarin for to urge someone, then the Mandarin for to not allow.
\end{quote}

\subsection*{You'll want to know: biāozhǔn}
\emph{biāozhǔn} — \textquotedblleft{}standard to measure by\textquotedblright{}.

\begin{quote}
\emph{wǒ juéde zhège biāozhǔn tài gāo le} — I think the standard is set too high.
\end{quote}

Its characters come later; for now it is a word for the ear.

\begin{cousinweb}
\emph{biāozhǔn} is built from \emph{biāo} and \zh{准} 'level', both inherited from Middle Chinese, \emph{pjiew} and \emph{tsywinX} (Guangyun, in Baxter's transcription), which glosses the first 'to raise' and 'the tip of a tree', and the second 'level'.
\end{cousinweb}

\begin{grammarlens}[title={What you've built}]
One more word for giving reasons.
\end{grammarlens}

\subsection*{Your turn}
\begin{itemize}
\raggedright
  \item \emph{Say it:} \emph{biāozhǔn}
  \item \emph{Say it:} \emph{biāozhǔn}, once more
  \item \emph{Say it:} \emph{wǒ juéde zhège biāozhǔn tài gāo le} — I think the standard is set too high.
  \item \emph{Recall:} say the Mandarin for to urge someone, then the Mandarin for to not allow, then say \emph{biāozhǔn} again
\end{itemize}

\begin{tcolorbox}[breakable,colback=teal!4,colframe=teal!35!black,title={Before you move on}]
What does biāozhǔn mean? (\textquotedblleft{}Standard to measure by\textquotedblright{}.) Say it once more.
\end{tcolorbox}
\section[pīzhǔn]{pīzhǔn — to give official approval; to approve a plan}

\label{lesson:ZH-C364-pizhun}



\begin{quote}
Before the new one: say the Mandarin for to not allow, then the Mandarin for a standard to measure by.
\end{quote}

\subsection*{You'll want to know: pīzhǔn}
\emph{pīzhǔn} — \textquotedblleft{}to give official approval; to approve a plan\textquotedblright{}.

\begin{quote}
\emph{xiàozhǎng pīzhǔn le wǒmen de jìhuà} — The head approved our plan.
\end{quote}

Its characters come later; for now it is a word for the ear.

\begin{cousinweb}
\emph{pīzhǔn} is built from \emph{pī} and \zh{准} 'level', both inherited from Middle Chinese, \emph{phej} and \emph{tsywinX} (Guangyun, in Baxter's transcription), which glosses the first 'to strike' and 'to show', and the second 'level'. How the pair came to mean approving is not settled.
\end{cousinweb}

\begin{grammarlens}[title={What you've built}]
One more word for giving reasons.
\end{grammarlens}

\subsection*{Your turn}
\begin{itemize}
\raggedright
  \item \emph{Say it:} \emph{pīzhǔn}
  \item \emph{Say it:} \emph{pīzhǔn}, once more
  \item \emph{Say it:} \emph{xiàozhǎng pīzhǔn le wǒmen de jìhuà} — The head approved our plan.
  \item \emph{Recall:} say the Mandarin for to not allow, then the Mandarin for a standard to measure by, then say \emph{pīzhǔn} again
\end{itemize}

\begin{tcolorbox}[breakable,colback=teal!4,colframe=teal!35!black,title={Before you move on}]
What does pīzhǔn mean? (\textquotedblleft{}To give official approval; to approve a plan\textquotedblright{}.) Say it once more.
\end{tcolorbox}
\section[tíxǐng]{tíxǐng — to remind}

\label{lesson:ZH-C364-tixing}



\begin{quote}
Before the new one: say the Mandarin for a standard to measure by, then the Mandarin for to give official approval.
\end{quote}

\subsection*{You'll want to know: tíxǐng}
\emph{tíxǐng} — \textquotedblleft{}to remind\textquotedblright{}.

\begin{quote}
\emph{qǐng tíxǐng wǒ míngtiān dài shū} — Please remind me to bring the book tomorrow.
\end{quote}

Its characters come later; for now it is a word for the ear.

\begin{cousinweb}
\emph{tíxǐng} is built from \emph{tí} and \emph{xǐng}, both inherited from Middle Chinese, \emph{dej} and \emph{sengX} (Guangyun, in Baxter's transcription), which glosses the first 'to carry by the hand' and the second 'to sober up'.
\end{cousinweb}

\begin{grammarlens}[title={What you've built}]
One more word for giving reasons.
\end{grammarlens}

\subsection*{Your turn}
\begin{itemize}
\raggedright
  \item \emph{Say it:} \emph{tíxǐng}
  \item \emph{Say it:} \emph{tíxǐng}, once more
  \item \emph{Say it:} \emph{qǐng tíxǐng wǒ míngtiān dài shū} — Please remind me to bring the book tomorrow.
  \item \emph{Recall:} say the Mandarin for a standard to measure by, then the Mandarin for to give official approval, then say \emph{tíxǐng} again
\end{itemize}

\begin{tcolorbox}[breakable,colback=teal!4,colframe=teal!35!black,title={Before you move on}]
What does tíxǐng mean? (\textquotedblleft{}To remind\textquotedblright{}.) Say it once more.
\end{tcolorbox}
